\chapter*{Referências}
\addcontentsline{toc}{chapter}{Referências}
\begin{singlespace}\setlength{\parindent}{0pt}\setlength{\parskip}{8pt}\raggedright

ABADIE, A. Using synthetic controls: feasibility, data requirements, and methodological aspects. \textbf{Journal of Economic Literature}, v. 59, n. 2, p. 391-425, 2021.

ABADIE, A.; DIAMOND, A.; HAINMUELLER, J. Synthetic control methods for comparative case studies: estimating the effect of California's tobacco control program. \textbf{Journal of the American Statistical Association}, v. 105, n. 490, p. 493-505, 2010.

ALEMAYEHU, F. K.; TVETERAAS, S. L. Long-run labour flexibility in
hospitality: a dynamic common correlated effects approach. \textbf{Tourism
Economics}, v. 26, n. 4, p. 704-718, 2020.

ANGRIST, J. D.; PISCHKE, J.-S. \textbf{Mostly harmless econometrics}: an empiricist's companion. Princeton: Princeton University Press, 2009.

ARELLANO, M.; BOND, S. Some tests of specification for panel data: Monte Carlo evidence and an application to employment equations. \textbf{The Review of Economic Studies}, v. 58, n. 2, p. 277-297, 1991.

ARELLANO, M.; BOVER, O. Another look at the instrumental variable estimation of error-components models. \textbf{Journal of Econometrics}, v. 68, n. 1, p. 29-51, 1995.

BAADE, R. A.; MATHESON, V. A. The quest for the cup: assessing the economic impact of the World Cup. \textbf{Regional Studies}, v. 38, n. 4, p. 343-354, 2004.

BAUMANN, R.; MATHESON, V. A. Mega-events and tourism: the case of Brazil. \textbf{Contemporary Economic Policy}, v. 36, n. 2, 2018.

BERTRAND, M.; DUFLO, E.; MULLAINATHAN, S. How much should we trust differences-in-differences estimates? \textbf{The Quarterly Journal of Economics}, v. 119, n. 1, p. 249-275, 2004.

BILLINGS, S. B.; HOLLADAY, J. S. Should cities go for the gold? The long-term impacts of hosting the Olympics. \textbf{Economic Inquiry}, v. 50, n. 3, p. 754-772, 2012.

BLUNDELL, R.; BOND, S. Initial conditions and moment restrictions in dynamic panel data models. \textbf{Journal of Econometrics}, v. 87, n. 1, p. 115-143, 1998.

BOTO-GARCÍA, D.; SANTANA-GALLEGO, M. Old questions, new methods: revisiting
the economic effects of hosting mega-sport events. \textbf{Tourism
Management}, v. 113, 105324, 2025.

BRASIL. Ministério do Trabalho e Emprego. \textbf{Portaria MTE nº 2.420, de
10 de julho de 2023}. Anexo XV: metodologia do Novo CAGED. Brasília, 2023.

BRONZINI, R.; MOCETTI, S.; MONGARDINI, M. \textbf{The economic effects of
big events}: evidence from the Great Jubilee 2000 in Rome. Roma: Banca
d'Italia, 2019. (Temi di discussione, n. 1208).

CARD, D.; KRUEGER, A. B. Minimum wages and employment: a case study of the fast-food industry in New Jersey and Pennsylvania. \textbf{American Economic Review}, v. 84, n. 4, p. 772-793, 1994.

CONLEY, T. G.; TABER, C. R. Inference with ``difference in differences'' with a small number of policy changes. \textbf{Review of Economics and Statistics}, v. 93, n. 1, p. 113-125, 2011.

CROMPTON, J. L. Economic impact analysis of sports facilities and events:
eleven sources of misapplication. \textbf{Journal of Sport Management}, v. 9,
n. 1, p. 14-35, 1995.

CUNNINGHAM, S. \textbf{Causal inference}: the mixtape. New Haven: Yale University Press, 2021.

DAVIS, S. J.; HALTIWANGER, J. Gross job creation, gross job destruction, and
employment reallocation. \textbf{The Quarterly Journal of Economics}, v. 107,
n. 3, p. 819-863, 1992.

DIX-CARNEIRO, R.; KOVAK, B. K. Trade liberalization and regional dynamics. \textbf{American Economic Review}, v. 107, n. 10, p. 2908-2946, 2017.

DOGRU, T.; McGINLEY, S.; SHARMA, A.; IŞIK, C.; HANKS, L. Employee turnover
dynamics in the hospitality industry vs. the overall economy.
\textbf{Tourism Management}, v. 99, 104783, 2023.

DOMINGUES, E. P.; BETARELLI JUNIOR, A. A.; MAGALHÃES, A. S. Quanto vale o
show? Impactos econômicos dos investimentos da Copa do Mundo 2014 no Brasil.
\textbf{Estudos Econômicos}, v. 41, n. 2, p. 409-439, 2011.

DWYER, L.; FORSYTH, P.; SPURR, R. Estimating the impacts of special events on an economy. \textbf{Journal of Travel Research}, v. 43, n. 4, p. 351-359, 2005.

FABER, B.; GAUBERT, C. Tourism and economic development: evidence from Mexico's coastline. \textbf{American Economic Review}, v. 109, n. 6, p. 2245-2293, 2019.

FEDDERSEN, A.; MAENNIG, W. Mega-events and sectoral employment: the case of
the 1996 Olympic Games. \textbf{Contemporary Economic Policy}, v. 31, n. 3,
p. 580-603, 2013.

FOURIE, J.; SANTANA-GALLEGO, M. The impact of mega-sport events on tourist arrivals. \textbf{Tourism Management}, v. 32, n. 6, p. 1364-1370, 2011.

HAGN, F.; MAENNIG, W. Large sport events and unemployment: the case of the
2006 soccer World Cup in Germany. \textbf{Applied Economics}, v. 41, n. 25,
p. 3295-3302, 2009.

HARVEY, D. From managerialism to entrepreneurialism: the transformation in
urban governance in late capitalism. \textbf{Geografiska Annaler: Series B,
Human Geography}, v. 71, n. 1, p. 3-17, 1989.

HICKS, J. R. \textbf{The theory of wages}. London: Macmillan, 1932.

HUNTINGTON-KLEIN, N. \textbf{The effect}: an introduction to research design and causality. Boca Raton: CRC Press, 2021.

IPEA. Substituição da captação dos dados do Caged pelo eSocial: implicações
para as estatísticas do emprego formal. \textbf{Boletim Mercado de Trabalho:
conjuntura e análise}, n. 69, 2020.

KOBIERECKI, M. M.; PIERZGALSKI, M. Sports mega-events and economic growth: a
synthetic control approach. \textbf{Journal of Sports Economics}, v. 23,
n. 5, p. 567-597, 2022.

KRAKOVER, S. Partitioning seasonal employment in the hospitality industry. \textbf{Tourism Management}, v. 21, n. 5, p. 461-471, 2000.

LEE, C.-K.; TAYLOR, T. Critical reflections on the economic impact assessment of a mega-event: the case of 2002 FIFA World Cup. \textbf{Tourism Management}, v. 26, n. 4, p. 595-603, 2005.

LEONI, V.; BOTO-GARCÍA, D. The effect of natural disasters on hotel demand, supply and labour markets: evidence from the La Palma volcano eruption. \textbf{Environmental and Resource Economics}, v. 86, n. 4, p. 755-780, 2023.

MATHESON, V. A. Economic multipliers and mega-event analysis.
\textbf{International Journal of Sport Finance}, v. 4, n. 1, p. 63-70, 2009.

MATTEI, T. S.; MELLO, G. R. Copa do mundo de futebol 2014: análise do
impacto da receita e volume de vendas do comércio, dos serviços e da produção
industrial nos estados brasileiros. \textbf{Revista Gestão \& Tecnologia},
v. 17, n. 1, p. 207-231, 2017.

MEURER, R.; LINS, H. N. The effects of the 2014 World Cup and the 2016 Olympic Games on Brazilian international travel receipts. \textbf{Tourism Economics}, v. 24, n. 4, p. 486-491, 2018.

MORTENSEN, D. T.; PISSARIDES, C. A. Job creation and job destruction in the
theory of unemployment. \textbf{The Review of Economic Studies}, v. 61, n. 3,
p. 397-415, 1994.

NEWEY, W. K.; WEST, K. D. A simple, positive semi-definite, heteroskedasticity and autocorrelation consistent covariance matrix. \textbf{Econometrica}, v. 55, n. 3, p. 703-708, 1987.

NICKELL, S. Biases in dynamic models with fixed effects. \textbf{Econometrica}, v. 49, n. 6, p. 1417-1426, 1981.

PREUSS, H. The conceptualisation and measurement of mega sport event
legacies. \textbf{Journal of Sport \& Tourism}, v. 12, n. 3-4, p. 207-228,
2007.

ROSE, A. K.; SPIEGEL, M. M. The Olympic effect. \textbf{The Economic Journal}, v. 121, n. 553, p. 652-677, 2011.

ROTH, J.; SANT'ANNA, P. H. C.; BILINSKI, A.; POE, J. What's trending in difference-in-differences? A synthesis of the recent econometrics literature. \textbf{Journal of Econometrics}, v. 235, n. 2, p. 2218-2244, 2023.

WAN, S. K.; SONG, H. Economic impact assessment of mega-events in the United
Kingdom and Brazil. \textbf{Journal of Hospitality \& Tourism Research},
v. 43, n. 7, p. 1044-1067, 2019.

\end{singlespace}
